\documentclass[10pt,a4paper]{amsart}
\usepackage{amsmath,amssymb,amsthm,graphicx}
\usepackage{verbatim}
\newenvironment{reference}{\comment}{\endcomment}
%\newenvironment{reference}{}{}
\newenvironment{slogan}{\comment}{\endcomment}
\newenvironment{history}{\comment}{\endcomment}
% For commutative diagrams we use Xy-pic
\usepackage[all]{xy}

% We use 2cell for 2-commutative diagrams.
\xyoption{2cell}
% Theorem environments.
%
\theoremstyle{plain}
\newtheorem{theorem}[subsection]{Theorem}
\newtheorem{proposition}[subsection]{Proposition}
\newtheorem{lemma}[subsection]{Lemma}

\theoremstyle{definition}
\newtheorem{definition}[subsection]{Definition}
\newtheorem{example}[subsection]{Example}
\newtheorem{exercise}[subsection]{Exercise}
\newtheorem{situation}[subsection]{Situation}

\theoremstyle{remark}
\newtheorem{remark}[subsection]{Remark}
\newtheorem{remarks}[subsection]{Remarks}

\numberwithin{equation}{subsection}

% Macros
%
\def\lim{\mathop{\mathrm{lim}}\nolimits}
\def\colim{\mathop{\mathrm{colim}}\nolimits}
\def\Spec{\mathop{\mathrm{Spec}}}
\def\Hom{\mathop{\mathrm{Hom}}\nolimits}
\def\Ext{\mathop{\mathrm{Ext}}\nolimits}
\def\SheafHom{\mathop{\mathcal{H}\!\mathit{om}}\nolimits}
\def\SheafExt{\mathop{\mathcal{E}\!\mathit{xt}}\nolimits}
\def\Sch{\mathit{Sch}}
\def\Mor{\mathop{\mathrm{Mor}}\nolimits}
\def\Ob{\mathop{\mathrm{Ob}}\nolimits}
\def\Sh{\mathop{\mathit{Sh}}\nolimits}
\def\NL{\mathop{N\!L}\nolimits}
\def\CH{\mathop{\mathrm{CH}}\nolimits}
\def\proetale{{pro\text{-}\acute{e}tale}}
\def\etale{{\acute{e}tale}}
\def\QCoh{\mathit{QCoh}}
\def\Ker{\mathop{\mathrm{Ker}}}
\def\Im{\mathop{\mathrm{Im}}}
\def\Coker{\mathop{\mathrm{Coker}}}
\def\Coim{\mathop{\mathrm{Coim}}}

% Boxtimes
%
\DeclareMathSymbol{\boxtimes}{\mathbin}{AMSa}{"02}

%
% Macros for moduli stacks/spaces
%
\def\QCohstack{\mathcal{QC}\!\mathit{oh}}
\def\Cohstack{\mathcal{C}\!\mathit{oh}}
\def\Spacesstack{\mathcal{S}\!\mathit{paces}}
\def\Quotfunctor{\mathrm{Quot}}
\def\Hilbfunctor{\mathrm{Hilb}}
\def\Curvesstack{\mathcal{C}\!\mathit{urves}}
\def\Polarizedstack{\mathcal{P}\!\mathit{olarized}}
\def\Complexesstack{\mathcal{C}\!\mathit{omplexes}}
% \Pic is the operator that assigns to X its picard group, usage \Pic(X)
% \Picardstack_{X/B} denotes the Picard stack of X over B
% \Picardfunctor_{X/B} denotes the Picard functor of X over B
\def\Pic{\mathop{\mathrm{Pic}}\nolimits}
\def\Picardstack{\mathcal{P}\!\mathit{ic}}
\def\Picardfunctor{\mathrm{Pic}}
\def\Deformationcategory{\mathcal{D}\!\mathit{ef}}
\usepackage{cleveref}
\crefname{originallemma}{original lemma}{original lemmas}
\Crefname{originallemma}{Original lemma}{Original lemmas}
\makeatletter
\newlabel{cohomology-lemma-projection-formula}{{01E8}{source}{}{sourcepointer-1}{}}
\newlabel{cohomology-lemma-projection-formula@cref}{{[originallemma][0][]01E8}{[1][0][]source}{}{}{}}
\makeatother
\crefname{originallemma}{original lemma}{original lemmas}
\Crefname{originallemma}{Original lemma}{Original lemmas}
\makeatletter
\newlabel{lemma-de-rham-complex-product}{{0FMA}{source}{}{sourcepointer-2}{}}
\newlabel{lemma-de-rham-complex-product@cref}{{[originallemma][0][]0FMA}{[1][0][]source}{}{}{}}
\makeatother
\crefname{originallemma}{original lemma}{original lemmas}
\Crefname{originallemma}{Original lemma}{Original lemmas}
\makeatletter
\newlabel{lemma-hodge-cohomology-projective-space}{{0FMI}{source}{}{sourcepointer-3}{}}
\newlabel{lemma-hodge-cohomology-projective-space@cref}{{[originallemma][0][]0FMI}{[1][0][]source}{}{}{}}
\makeatother
\crefname{originallemma}{original lemma}{original lemmas}
\Crefname{originallemma}{Original lemma}{Original lemmas}
\makeatletter
\newlabel{lemma-push-omega-a}{{0FUH}{source}{}{sourcepointer-4}{}}
\newlabel{lemma-push-omega-a@cref}{{[originallemma][0][]0FUH}{[1][0][]source}{}{}{}}
\makeatother
\crefname{originallemma}{original lemma}{original lemmas}
\Crefname{originallemma}{Original lemma}{Original lemmas}
\makeatletter
\newlabel{lemma-the-complex-for-L-star}{{0FUF}{source}{}{sourcepointer-5}{}}
\newlabel{lemma-the-complex-for-L-star@cref}{{[originallemma][0][]0FUF}{[1][0][]source}{}{}{}}
\makeatother
\crefname{originallemma}{original lemma}{original lemmas}
\Crefname{originallemma}{Original lemma}{Original lemmas}
\makeatletter
\newlabel{lemma-twisted-hodge-cohomology-projective-space}{{0FUK}{source}{}{sourcepointer-6}{}}
\newlabel{lemma-twisted-hodge-cohomology-projective-space@cref}{{[originallemma][0][]0FUK}{[1][0][]source}{}{}{}}
\makeatother
\crefname{originalsection}{original section}{original sections}
\Crefname{originalsection}{Original section}{Original sections}
\makeatletter
\newlabel{section-line-bundle}{{0FU8}{source}{}{sourcepointer-7}{}}
\newlabel{section-line-bundle@cref}{{[originalsection][0][]0FU8}{[1][0][]source}{}{}{}}
\makeatother
\crefname{originalsection}{original section}{original sections}
\Crefname{originalsection}{Original section}{Original sections}
\makeatletter
\newlabel{section-projective-space}{{0FMG}{source}{}{sourcepointer-8}{}}
\newlabel{section-projective-space@cref}{{[originalsection][0][]0FMG}{[1][0][]source}{}{}{}}
\makeatother
\title{Comparison of the differential forms on the blowup, the exceptional divisor and the punctured total space}
\author{The Stacks Project Authors}
\date{stacks-project:derham.tex:local-snapshot-20261012; GNU FDL 1.2 or later}
\begin{document}
\maketitle
\noindent\textit{Editorial note.} All delivered mathematical text is the author's own, copied from the pinned Stacks Project chapter file derham.tex; no statement or proof step was rewritten, reordered or removed.\par\medskip
\noindent\textit{Editorial note.} Preamble provenance: the pinned chapter file opens with \textbackslash{}input\{preamble\} and carries no preamble text of its own, so the project's single Stacks-wide preamble staged at sources/stacks-derham/preamble.tex is excerpted verbatim as three spans, lines 8-12 (verbatim package and the project's own reference/slogan/history display environments), lines 14-18 (Xy-pic setup) and lines 154-217 (theorem declarations and mathematical macros). The wrapper class and the amsmath/amssymb/amsthm/graphicx base packages are supplied by the assembler. No macro was invented or redefined.\par\medskip
\noindent\textit{Editorial note.} Same-chapter pointers are located at their \textbackslash{}label lines inside the pinned chapter file; cross-chapter pointers are located in the project's own tag list staged in the snapshot (sources/stacks-chow/tags.txt), at the verbatim lines carrying their canonical tags.\par\medskip
\noindent\textit{Allowed standard prerequisites (editorial).}\begin{itemize}
\item The projection formula. Cohomology, label cohomology-lemma-projection-formula (invoked at proof line 2877). It rewrites \$p\_*\$ of the product decomposition of \$\textbackslash{}Omega\textasciicircum{}a\_\{P/S\}\$ as a direct sum of tensor products with \$\textbackslash{}operatorname\{pr\}\_\{1,*\}\textbackslash{}operatorname\{pr\}\_2\textasciicircum{}*\textbackslash{}Omega\textasciicircum{}s\$.
\item The computation of the Hodge cohomology of projective space, and the section on projective space. Same chapter derham.tex, Section \textbackslash{}ref\{section-projective-space\} (invoked at proof line 2884) and label lemma-hodge-cohomology-projective-space (invoked at proof line 2886). They show that \$\textbackslash{}operatorname\{pr\}\_\{1,*\}\textbackslash{}operatorname\{pr\}\_2\textasciicircum{}*\textbackslash{}Omega\textasciicircum{}s\_\{\textbackslash{}mathbf\{P\}\textasciicircum{}\{n-1\}\_S/S\}\$ vanishes except for \$s = 0\$, which settles part (2).
\item The direct sum decomposition of the pushforward along the punctured line bundle and the exact sequences for its graded pieces. Same chapter derham.tex, Section \textbackslash{}ref\{section-line-bundle\} (invoked at proof line 2893) and labels lemma-push-omega-a (invoked at proof lines 2894 and 2900) and lemma-the-complex-for-L-star (invoked at proof line 2900). They give \$\textbackslash{}pi\_*\textbackslash{}Omega\textasciicircum{}a\_\{L/S\} = \textbackslash{}Omega\textasciicircum{}a\_\{P/S\} \textbackslash{}oplus \textbackslash{}bigoplus\_\{k\textbackslash{}geq 1\}\textbackslash{}Omega\textasciicircum{}a\_\{L/S,k\}\$ and the short exact sequences exhibiting \$\textbackslash{}Omega\textasciicircum{}a\_\{L/S,k\}\$ as an extension of twisted pieces, which is how part (3) is reduced.
\item The product decomposition of differential forms and the twisted Hodge cohomology of projective space. Same chapter derham.tex, label lemma-de-rham-complex-product (invoked at proof line 2913) and label lemma-twisted-hodge-cohomology-projective-space (invoked at proof line 2918). The first decomposes \$\textbackslash{}Omega\textasciicircum{}a\_\{P/S\}\$ over \$P = Z \textbackslash{}times\_S \textbackslash{}mathbf\{P\}\textasciicircum{}\{n-1\}\_S\$, the second kills the higher cohomology of the twisted pieces \$\textbackslash{}mathcal\{O\}\_Z \textbackslash{}boxtimes \textbackslash{}Omega\textasciicircum{}j(k)\$ for \$k > 0\$, completing part (3).
\end{itemize}\medskip
\ifdefined\hypertarget\hypertarget{sourcepointer-1}{}\fi
\noindent\textit{Source reference.} Original lemma 01E8: sources/stacks-chow/tags.txt, source lines 1825--1825. Navigation pointer for \textbackslash{}ref\{cohomology-lemma-projection-formula\} used inside the counted span (line 2877); the argument that consumes it is retained in the span.\par\medskip
\ifdefined\hypertarget\hypertarget{sourcepointer-2}{}\fi
\noindent\textit{Source reference.} Original lemma 0FMA: sources/stacks-derham/pinned-original-source.tex, source lines 871--871. Navigation pointer for \textbackslash{}ref\{lemma-de-rham-complex-product\} used inside the counted span (line 2913); the argument that consumes it is retained in the span.\par\medskip
\ifdefined\hypertarget\hypertarget{sourcepointer-3}{}\fi
\noindent\textit{Source reference.} Original lemma 0FMI: sources/stacks-derham/pinned-original-source.tex, source lines 1705--1705. Navigation pointer for \textbackslash{}ref\{lemma-hodge-cohomology-projective-space\} used inside the counted span (line 2886); the argument that consumes it is retained in the span.\par\medskip
\ifdefined\hypertarget\hypertarget{sourcepointer-4}{}\fi
\noindent\textit{Source reference.} Original lemma 0FUH: sources/stacks-derham/pinned-original-source.tex, source lines 1509--1509. Navigation pointer for \textbackslash{}ref\{lemma-push-omega-a\} used inside the counted span (line 2894); the argument that consumes it is retained in the span.\par\medskip
\ifdefined\hypertarget\hypertarget{sourcepointer-5}{}\fi
\noindent\textit{Source reference.} Original lemma 0FUF: sources/stacks-derham/pinned-original-source.tex, source lines 1348--1348. Navigation pointer for \textbackslash{}ref\{lemma-the-complex-for-L-star\} used inside the counted span (line 2900); the argument that consumes it is retained in the span.\par\medskip
\ifdefined\hypertarget\hypertarget{sourcepointer-6}{}\fi
\noindent\textit{Source reference.} Original lemma 0FUK: sources/stacks-derham/pinned-original-source.tex, source lines 1644--1644. Navigation pointer for \textbackslash{}ref\{lemma-twisted-hodge-cohomology-projective-space\} used inside the counted span (line 2918); the argument that consumes it is retained in the span.\par\medskip
\ifdefined\hypertarget\hypertarget{sourcepointer-7}{}\fi
\noindent\textit{Source reference.} Original section 0FU8: sources/stacks-derham/pinned-original-source.tex, source lines 1250--1250. Navigation pointer for \textbackslash{}ref\{section-line-bundle\} used inside the counted span (line 2893); the argument that consumes it is retained in the span.\par\medskip
\ifdefined\hypertarget\hypertarget{sourcepointer-8}{}\fi
\noindent\textit{Source reference.} Original section 0FMG: sources/stacks-derham/pinned-original-source.tex, source lines 1576--1576. Navigation pointer for \textbackslash{}ref\{section-projective-space\} used inside the counted span (line 2884); the argument that consumes it is retained in the span.\par\medskip
\begin{lemma}
\label{lemma-comparison}
For $a \geq 0$ we have
\begin{enumerate}
\item the map
$\Omega^a_{X/S} \to b_*\Omega^a_{L/S}$ is an isomorphism,
\item the map $\Omega^a_{Z/S} \to p_*\Omega^a_{P/S}$ is an isomorphism,
and
\item the map $Rb_*\Omega^a_{L/S} \to i_*Rp_*\Omega^a_{P/S}$ is an isomorphism
on cohomology sheaves in degree $\geq 1$.
\end{enumerate}
\end{lemma}

\begin{proof}
Let us first prove part (2). Since
$P = Z \times_S \mathbf{P}^{n - 1}_S$
we see that
$$
\Omega^a_{P/S} = \bigoplus\nolimits_{a = r + s}
\text{pr}_1^*\Omega^r_{Z/S} \otimes
\text{pr}_2^*\Omega^s_{\mathbf{P}^{n - 1}_S/S}
$$
Recalling that $p = \text{pr}_1$ by the projection formula
(Cohomology, Lemma \ref{cohomology-lemma-projection-formula})
we obtain
$$
p_*\Omega^a_{P/S} = \bigoplus\nolimits_{a = r + s}
\Omega^r_{Z/S} \otimes
\text{pr}_{1, *}\text{pr}_2^*\Omega^s_{\mathbf{P}^{n - 1}_S/S}
$$
By the calculations in Section \ref{section-projective-space}
and in particular in
the proof of Lemma \ref{lemma-hodge-cohomology-projective-space}
we have $\text{pr}_{1, *}\text{pr}_2^*\Omega^s_{\mathbf{P}^{n - 1}_S/S} = 0$
except if $s = 0$ in which case we get
$\text{pr}_{1, *}\mathcal{O}_P = \mathcal{O}_Z$.
This proves (2).

\medskip\noindent
By the material in Section \ref{section-line-bundle} and in particular
Lemma \ref{lemma-push-omega-a} we have
$\pi_*\Omega^a_{L/S} = \Omega^a_{P/S} \oplus
\bigoplus_{k \geq 1} \Omega^a_{L/S, k}$.
Since the composition $\pi \circ 0$ in the diagram above
is the identity morphism on $P$ to prove part (3) it suffices to show that
$\Omega^a_{L/S, k}$ has vanishing higher cohomology for $k > 0$.
By Lemmas \ref{lemma-the-complex-for-L-star} and \ref{lemma-push-omega-a}
there are short exact sequences
$$
0 \to \Omega^a_{P/S} \otimes \mathcal{O}_P(k)
\to \Omega^a_{L/S, k} \to
\Omega^{a - 1}_{P/S} \otimes \mathcal{O}_P(k) \to 0
$$
where $\Omega^{a - 1}_{P/S} = 0$ if $a = 0$. Since
$P = Z \times_S \mathbf{P}^{n - 1}_S$ we have
$$
\Omega^a_{P/S} = \bigoplus\nolimits_{i + j = a}
\Omega^i_{Z/S} \boxtimes \Omega^j_{\mathbf{P}^{n - 1}_S/S}
$$
by Lemma \ref{lemma-de-rham-complex-product}.
Since $\Omega^i_{Z/S}$ is free of finite rank
we see that it suffices to show that the higher cohomology of
$\mathcal{O}_Z \boxtimes \Omega^j_{\mathbf{P}^{n - 1}_S/S}(k)$
is zero for $k > 0$. This follows from
Lemma \ref{lemma-twisted-hodge-cohomology-projective-space}
applied to $P = Z \times_S \mathbf{P}^{n - 1}_S = \mathbf{P}^{n - 1}_Z$
and the proof of (3) is complete.

\medskip\noindent
We still have to prove (1). If $n = 1$, then we are blowing
up an effective Cartier divisor and $b$ is an isomorphism
and we have (1). If $n > 1$, then the composition
$$
\Gamma(X, \Omega^a_{X/S})
\to
\Gamma(L, \Omega^a_{L/S})
\to
\Gamma(L \setminus E, \Omega^a_{L/S})
=
\Gamma(X \setminus Z, \Omega^a_{X/S})
$$
is an isomorphism as $\Omega^a_{X/S}$ is finite free
(small detail omitted). Thus the only way (1) can fail is if
there are nonzero elements of $\Gamma(L, \Omega^a_{L/S})$ which vanish
outside of $E = 0(P)$. Since $L$ is a line bundle over $P$
with zero section $0 : P \to L$, it suffices to show that
on a line bundle there are no nonzero sections of a sheaf
of differentials which vanish identically outside the zero section.
The reader sees this is true either (preferably) by a local calculation
or by using that $\Omega_{L/S, k} \subset \Omega_{L^\star/S, k}$
(see references above).
\end{proof}

\end{document}
